\chapter{Metric Spaces and Completeness}

\begin{center}\itshape
Iterate $x\mapsto\cos x$ from any starting point and the values settle on $0.739085\ldots$; iterate two maps that shrink
the line by a factor $3$ and the Cantor set appears, whatever shape you start from. Why must such iterations
converge, and why always to the same limit?
\end{center}

Take a pocket calculator, set it to radians, type any number and press the cosine key, again and again. After a few
dozen presses the display freezes on $0.7390851\ldots$, whatever number you started from. A professor of French named
Dottie is said to have noticed this while idly pressing keys, and the number now carries her name. Why does it happen?
Because $\cos$ brings points closer together, and a map that brings points closer together, on a space without holes,
has exactly one point that it does not move. Repeating the map walks you there.

To say this properly one needs only a notion of distance, and Maurice Fr\'echet realised in his thesis of 1906 that a
distance is \emph{all} one needs: functions, curves and sequences can be points of a space as soon as one says how
far apart they are. Felix Hausdorff named these spaces \emph{metric spaces} in 1914. The key property is
\emph{completeness}: Cauchy sequences converge, there are no holes. Chapter~1 completed $\Q$ into $\R$; the same
construction completes every metric space, and completing $\Q$ for the strange $p$-adic distance of Tome~1 produces
another universe of numbers, $\Q_p$.

The chapter's hero is the fixed point theorem of Stefan Banach (1922). It solves equations, proves that differential
equations have solutions, makes Newton's method converge, ranks the pages of the web, and draws fractals: the Cantor
set and the Sierpi\'nski triangle are the fixed points of maps acting on \emph{shapes}. Along the way, $\exp$, $\cos$,
$\sin$ and $\pi$, used since Tome~0, are finally defined, by power series, and $\pi$ appears not from a circle but as
twice the first zero of the cosine.

\section{Distances}

\begin{definition}
A \emph{metric} on a set $X$ is a map $d:X\times X\to[0,+\infty[$ with $d(x,y)=0$ iff $x=y$, $d(x,y)=d(y,x)$, and the triangle inequality
$d(x,z)\le d(x,y)+d(y,z)$. A \emph{norm} on a real or complex vector space $E$ is a map $\|\cdot\|:E\to[0,+\infty[$ with $\|x\|=0$ iff $x=0$,
$\|\lambda x\|=|\lambda|\|x\|$ and $\|x+y\|\le\|x\|+\|y\|$; then $d(x,y)=\|x-y\|$ is a metric.
\end{definition}

\begin{example}\label{ex:metrics}
\begin{enumerate}[label=(\alph*)]
\item On $\R^n$: $\|x\|_1=\sum|x_i|$, $\|x\|_2=\bigl(\sum x_i^2\bigr)^{1/2}$ and $\|x\|_\infty=\max|x_i|$. The triangle inequality for $\|\cdot\|_2$ is Minkowski's, a
consequence of Cauchy--Schwarz: $\|x+y\|_2^2=\|x\|_2^2+2\langle x,y\rangle+\|y\|_2^2\le(\|x\|_2+\|y\|_2)^2$. They satisfy
$\|x\|_\infty\le\|x\|_2\le\|x\|_1\le n\|x\|_\infty$.
\item On the set $B(S)$ of bounded functions $S\to\R$ (or $\C$), the \emph{sup norm} $\|f\|_\infty=\sup_{s\in S}|f(s)|$.
\item On any set, the \emph{discrete metric}: $d(x,y)=1$ if $x\ne y$. On words of length $n$, the \emph{Hamming distance}: the number of
positions where two words differ, the distance of error-correcting codes.
\item On $\Q$, for a prime $p$, the \emph{$p$-adic absolute value} $|x|_p=p^{-v_p(x)}$, where $v_p(x)$ is the exponent of $p$ in $x\ne0$, and $|0|_p=0$
(Tome~1, Chapter~12). It satisfies $|xy|_p=|x|_p|y|_p$ and the \emph{ultrametric} inequality $|x+y|_p\le\max(|x|_p,|y|_p)$, stronger than the
triangle inequality, because $v_p(x+y)\ge\min(v_p(x),v_p(y))$. So $d(x,y)=|x-y|_p$ is a metric, in
which $p^n\to0$.
\end{enumerate}
\end{example}

\begin{figure}[ht]
\centering
\fig{ch02-balls}
\caption{The unit balls of $\|\cdot\|_1$ (a diamond), $\|\cdot\|_2$ (a disc) and $\|\cdot\|_\infty$ (a square) in $\R^2$, nested as the inequalities of
Example~\ref{ex:metrics}(a) say. Each norm gives the same convergent sequences.}
\label{fig:balls}
\end{figure}

A subset $Y$ of a metric space is a metric space with the restricted metric. The \emph{open ball} is $B(x,r)=\{y:d(x,y)<r\}$, the
\emph{closed ball} $\bar B(x,r)=\{y:d(x,y)\le r\}$, and the \emph{diameter} of $A$ is $\operatorname{diam}A=\sup_{x,y\in A}d(x,y)$. The \emph{distance} from $x$ to a
non-empty $A$ is $d(x,A)=\inf_{a\in A}d(x,a)$; it is $1$-Lipschitz in $x$: $d(x,A)\le d(x,a)\le d(x,y)+d(y,a)$, and take the infimum over $a$.

\section{Open sets, convergence and continuity}

\begin{definition}
A set $U\subseteq X$ is \emph{open} if for every $x\in U$ there is $r>0$ with $B(x,r)\subseteq U$; it is \emph{closed} if $X\setminus U$ is open. The
\emph{interior} $\mathring A$ is the union of the open sets inside $A$, the \emph{closure} $\bar A$ the intersection of the closed sets containing $A$,
and $A$ is \emph{dense} if $\bar A=X$. A sequence $x_n\to x$ if $d(x_n,x)\to0$.
\end{definition}

An open ball is open: if $y\in B(x,r)$, then $B(y,r-d(x,y))\subseteq B(x,r)$ by the triangle inequality. Unions of open sets and finite
intersections of open sets are open (take the smallest radius), so intersections of closed sets and finite unions are
closed. A limit is unique: $d(x,y)\le d(x,x_n)+d(x_n,y)$.

\begin{proposition}\label{prop:closure}
\begin{enumerate}[label=(\alph*)]
\item $x\in\bar A$ iff $d(x,A)=0$ iff some sequence of $A$ converges to $x$.
\item $A$ is closed iff it contains the limits of its convergent sequences.
\item For $f:X\to Y$ between metric spaces, the following are equivalent: (i) for every $x$ and $\varepsilon>0$ there is $\delta>0$ with
$d(f(x),f(y))<\varepsilon$ when $d(x,y)<\delta$ (\emph{continuity}); (ii) $x_n\to x$ implies $f(x_n)\to f(x)$; (iii) $f^{-1}(V)$ is open for every
open $V\subseteq Y$.
\end{enumerate}
\end{proposition}
\begin{proof}
(a) If $d(x,A)>0$, the ball $B(x,d(x,A))$ is open and misses $A$, so its complement is a closed set containing $A$ and not $x$: $x\notin\bar A$.
If $d(x,A)=0$, pick $a_n\in A$ with $d(x,a_n)<\frac1n$: then $a_n\to x$. If $a_n\to x$ with $a_n\in A$, and $F\supseteq A$ is closed with $x\notin F$, some ball
$B(x,r)$ misses $F$ and hence all $a_n$, which is absurd for large $n$; so $x\in\bar A$. (b) $A$ is closed iff $A=\bar A$; apply (a).
(c) (i)$\Rightarrow$(ii) is direct. (ii)$\Rightarrow$(iii): if $f^{-1}(V)$ is not open, there are $x\in f^{-1}(V)$ and $x_n\to x$ with $f(x_n)\notin V$; then
$f(x_n)\to f(x)\in V$, but $V$ contains a ball around $f(x)$, which $f(x_n)$ must eventually enter. (iii)$\Rightarrow$(i): $f^{-1}(B(f(x),\varepsilon))$ is open and
contains $x$, so it contains a ball $B(x,\delta)$.
\end{proof}

By (c), continuity only depends on the open sets, not on the metric: this is the starting point of Chapter~3. Two
metrics are \emph{equivalent} if they have the same open sets, \emph{Lipschitz equivalent} if $c\,d\le d'\le C\,d$ for constants
$0<c\le C$. The three norms of Example~\ref{ex:metrics}(a) are Lipschitz equivalent, and Chapter~4 proves that all norms on $\R^n$ are.

\begin{definition}
A map $f:X\to Y$ is \emph{uniformly continuous} if the $\delta$ in (i) can be chosen independently of $x$, and \emph{$k$-Lipschitz} if
$d(f(x),f(y))\le k\,d(x,y)$ for all $x,y$. A \emph{contraction} is a $k$-Lipschitz map with $k<1$. An \emph{isometry} preserves distances.
\end{definition}

Lipschitz maps are uniformly continuous; $x\mapsto x^2$ on $\R$ is continuous but not uniformly ($(n+\frac1n)^2-n^2>2$), and $\sqrt x$
on $[0,1]$ is uniformly continuous but not Lipschitz.

\section{Completeness}

\begin{definition}
A sequence $(x_n)$ is \emph{Cauchy} if for every $\varepsilon>0$ there is $N$ with $d(x_n,x_m)<\varepsilon$ for $n,m\ge N$. The space is \emph{complete} if
every Cauchy sequence converges. A complete normed space is a \emph{Banach space}.
\end{definition}

A convergent sequence is Cauchy, and a Cauchy sequence with a convergent subsequence converges to the same limit
($d(x_n,x)\le d(x_n,x_{n_k})+d(x_{n_k},x)$). $\R$ is complete (Theorem~\ref{thm:reals}), and $\Q$ and $]0,1]$ are not: $\frac1n$ is Cauchy in
$]0,1]$ with no limit there. Completeness, unlike openness, is not a topological property: $]0,1]$ and $[1,+\infty[$ are
homeomorphic by $x\mapsto\frac1x$, and only the second is complete.

\begin{proposition}\label{prop:complete}
\begin{enumerate}[label=(\alph*)]
\item A closed subset of a complete space is complete; a complete subset of any metric space is closed.
\item $\R^n$ and $\C^n$ are complete for $\|\cdot\|_1$, $\|\cdot\|_2$ and $\|\cdot\|_\infty$.
\item $B(S)$ with the sup norm is a Banach space, and if $S$ is a metric space, the bounded continuous functions form a
closed subspace $C_b(S)$, hence a Banach space. In particular a uniform limit of continuous functions is continuous.
\item In a Banach space, an absolutely convergent series converges: if $\sum\|u_k\|<\infty$, then $\sum u_k$ converges and
$\|\sum u_k\|\le\sum\|u_k\|$.
\end{enumerate}
\end{proposition}
\begin{proof}
(a) A Cauchy sequence of a closed $F$ converges in $X$, and its limit is in $F$ by Proposition~\ref{prop:closure}(b). If $A$ is complete and
$a_n\to x$ with $a_n\in A$, then $(a_n)$ is Cauchy, so it converges in $A$, and by uniqueness of limits $x\in A$.
(b) By the inequalities between the norms, a sequence is Cauchy (or converges) for one norm iff it is for all, and for $\|\cdot\|_\infty$
this happens iff each coordinate is Cauchy (or converges) in $\R$. Identify $\C^n$ with $\R^{2n}$.
(c) Let $(f_n)$ be Cauchy in $B(S)$. For each $s$, $|f_n(s)-f_m(s)|\le\|f_n-f_m\|_\infty$, so $(f_n(s))$ is Cauchy and has a limit $f(s)$. Given $\varepsilon$,
take $N$ with $\|f_n-f_m\|_\infty\le\varepsilon$ for $n,m\ge N$, and let $m\to\infty$: $|f_n(s)-f(s)|\le\varepsilon$ for all $s$ and $n\ge N$. So $f$ is bounded and
$f_n\to f$ uniformly. If moreover each $f_n$ is continuous, fix $s$ and $\varepsilon$, take $n$ with $\|f_n-f\|_\infty<\varepsilon$ and $\delta$ with
$|f_n(t)-f_n(s)|<\varepsilon$ for $d(s,t)<\delta$; then $|f(t)-f(s)|<3\varepsilon$. So $C_b(S)$ is closed in $B(S)$, and complete by (a).
(d) The partial sums satisfy $\|\sum_{k=m}^nu_k\|\le\sum_{k=m}^n\|u_k\|$, which is small for large $m$: they are Cauchy.
\end{proof}

\section{The completion}

\begin{lemma}[extension]\label{lem:extension}
Let $D$ be dense in $X$, $Y$ complete, and $f:D\to Y$ uniformly continuous. Then $f$ has a unique continuous extension
$\bar f:X\to Y$, and $\bar f$ is uniformly continuous. If $f$ is an isometry, so is $\bar f$.
\end{lemma}
\begin{proof}
For $x\in X$ take $d_n\in D$ with $d_n\to x$. Then $(d_n)$ is Cauchy, so $(f(d_n))$ is Cauchy by uniform continuity, and converges in $Y$.
The limit does not depend on the sequence: interleaving two sequences converging to $x$ gives one sequence converging to $x$,
whose image converges, so the two limits agree. Call it $\bar f(x)$. Given $\varepsilon$, let $\delta$ be such that $d(a,b)<\delta$ implies
$d(f(a),f(b))<\varepsilon$ on $D$. If $d(x,y)<\delta$, approximating sequences $a_n\to x$, $b_n\to y$ eventually satisfy $d(a_n,b_n)<\delta$, so
$d(\bar f(x),\bar f(y))\le\varepsilon$: $\bar f$ is uniformly continuous. Uniqueness: two continuous extensions agree on $D$, hence on $X$ by
Proposition~\ref{prop:closure}(c)(ii). Distances pass to the limit, which gives the isometry statement.
\end{proof}

\begin{theorem}[completion]\label{thm:completion}
Every metric space $X$ has a \emph{completion}: a complete metric space $\hat X$ with an isometry $\iota:X\to\hat X$ whose image is dense. If
$(\hat X',\iota')$ is another one, there is a unique isometric bijection $\phi:\hat X\to\hat X'$ with $\phi\circ\iota=\iota'$.
\end{theorem}
\begin{proof}
\emph{Existence} (Kuratowski's embedding). We may assume $X\ne\varnothing$; fix $a\in X$. For $x\in X$ let $\iota(x)$ be the function $f_x(s)=d(x,s)-d(a,s)$ on $X$.
It is bounded, by $d(x,a)$, by the triangle inequality, and continuous, so $f_x\in C_b(X)$. Moreover
$\|f_x-f_y\|_\infty=\sup_s|d(x,s)-d(y,s)|=d(x,y)$, the supremum being attained at $s=y$. So $\iota$ is an isometry into the Banach space $C_b(X)$
(Proposition~\ref{prop:complete}(c)). Let $\hat X$ be the closure of $\iota(X)$; it is complete by Proposition~\ref{prop:complete}(a), and $\iota(X)$ is dense in it.
\emph{Uniqueness.} $\iota'\circ\iota^{-1}:\iota(X)\to\hat X'$ is an isometry from a dense subset into a complete space; by Lemma~\ref{lem:extension} it extends to an isometry
$\phi:\hat X\to\hat X'$. Its image is complete, hence closed, and contains the dense set $\iota'(X)$: $\phi$ is onto. Any $\phi$ with $\phi\circ\iota=\iota'$ is continuous and
is determined on the dense set $\iota(X)$.
\end{proof}

This proof cannot construct $\R$, since $C_b(X)$ is made of real-valued functions; that is why Chapter~1 used Cauchy
sequences. The same construction works here (classes of Cauchy sequences of $X$, with $d([x],[y])=\lim d(x_n,y_n)$), and gives a
second proof. When $X$ has more structure, the completion inherits it: the completion of a normed space is a Banach space
(the operations are uniformly continuous on bounded sets and extend), and Tome~7 finds $L^1$ and $L^2$ as completions of the
continuous functions for integral norms (Exercise~\ref{ex:L1incomplete}).

\section{The \texorpdfstring{$p$}{p}-adic numbers as a completion}

In 1897 Kurt Hensel had an idea that looked like a joke: treat a number divisible by a high power of $p$ as \emph{small}.
For the $p$-adic distance of Tome~1, $p^n\to0$, the series $1+p+p^2+\cdots$ converges (to $\frac1{1-p}$), and integers can be
very close to each other without being close at all on the real line. In 1916 Alexander Ostrowski proved that, up to
equivalence, the usual absolute value and the $p$-adic ones are the \emph{only} absolute values on $\Q$. So completing $\Q$
gives $\R$ and the fields $\Q_p$, and nothing else: number theory has one real place and one place for each prime, a
symmetry that Tome~12 exploits.


\begin{theorem}\label{thm:Qp}
Let $\Q_p$ be the completion of $(\Q,|\cdot|_p)$. Addition and multiplication extend from $\Q$, and $\Q_p$ is a field; $|\cdot|_p$ extends to an
ultrametric absolute value on $\Q_p$, with values in $\{p^k:k\in\Z\}\cup\{0\}$. The closed unit ball $\Z_p=\{x:|x|_p\le1\}$ is a subring, it is the
closure of $\Z$, and $x\mapsto(x\bmod p^n)_{n\ge1}$ is a ring isomorphism from $\Z_p$ onto the ring of compatible sequences that Tome~1,
Chapter~12 calls $\Z_p$. Every $x\in\Z_p$ is uniquely a convergent sum $\sum_{k\ge0}a_kp^k$ with digits $a_k\in\{0,\dots,p-1\}$.
\end{theorem}
\begin{proof}
Write $|\cdot|$ for $|\cdot|_p$. We use the Cauchy-sequence model of the completion: elements of $\Q_p$ are classes of Cauchy sequences of $\Q$,
and $\Q$ sits inside as constant sequences. Parts (a) and (b) of Lemma~\ref{lem:cauchybasic} of Chapter~1, and the bound $|x_n|\ge\delta$ for large $n$ in part (c),
hold word for word with $|\cdot|$ in place of the usual absolute value (their proofs used only the triangle inequality and
multiplicativity), so the classes form a field, exactly as in Theorem~\ref{thm:reals}. If $(x_n)$ is Cauchy and not null, then $|x_n|$ is eventually constant: $|x_n|\ge\delta$ for large $n$,
and once $|x_n-x_m|<\delta$, the ultrametric inequality gives $|x_n|=|x_m|$ (if $|a|>|b|$ then $|a+b|=|a|$, since $|a|\le\max(|a+b|,|b|)$). So
$|x|=\lim|x_n|$ is well defined, multiplicative and ultrametric (pass to the limit), with values in $p^\Z\cup\{0\}$, and it induces the
metric of the completion.

$\Z_p$ is a subring by the ultrametric inequality and multiplicativity. It is closed (as a closed ball) and contains $\Z$. If
$|x|\le1$, pick $q\in\Q$ with $|x-q|<p^{-n}$; then $|q|\le1$, so $q=\frac ab$ with $p\nmid b$; choose $c\in\Z$ with $bc\equiv1\pmod{p^n}$, and then
$|q-ac|=|a|\,|1-bc|/|b|\le p^{-n}$. So $x$ is within $p^{-n}$ of the integer $ac$: $\Z$ is dense in $\Z_p$.

For $x\in\Z_p$ and $n\ge1$, let $x\bmod p^n$ be the class modulo $p^n$ of any integer $m$ with $|x-m|\le p^{-n}$; two such integers differ by a
multiple of $p^n$, so this is well defined, and it is compatible with reduction from $p^{n+1}$ to $p^n$. The map is a ring morphism
(approximate $x$ and $y$ by integers $m$, $m'$; then $m+m'$ and $mm'$ approximate $x+y$ and $xy$ to within $p^{-n}$). It is injective: if all
the reductions vanish, $|x|\le p^{-n}$ for every $n$. It is onto: given compatible $(c_n)$, choose integers $m_n$ in the classes $c_n$; then
$|m_{n+1}-m_n|\le p^{-n}$, so $(m_n)$ is Cauchy, and its limit reduces to $c_n$ for each $n$.

Digits: let $a_0\in\{0,\dots,p-1\}$ be the representative of $x\bmod p$; then $|x-a_0|\le p^{-1}$, so $x-a_0=py$ with $|y|\le1$, and we continue with $y$. By
induction $|x-\sum_{k<n}a_kp^k|\le p^{-n}$, so the series converges to $x$. Uniqueness: if two digit sequences first differ at index $j$, the
difference of the sums is $p^j$ times $(a_j-a'_j)+p(\cdots)$, of absolute value exactly $p^{-j}$.
\end{proof}

So the $p$-adic integers of Tome~1, defined there by congruences, are the $p$-adic limits of ordinary integers. For example,
$\frac1{1-p}=\sum_{k\ge0}p^k$ in $\Z_p$, and $-1=\sum_{k\ge0}(p-1)p^k$: in $\Z_2$, $-1=\ldots1111$, the two's complement of computers.


\section{Banach's fixed point theorem}

In the summer of 1916, the mathematician Hugo Steinhaus was walking in the Planty park in Krak\'ow when he overheard the
words ``Lebesgue integral'' from a bench. He stopped and found two young men talking mathematics; one of them was Stefan
Banach, who had no degree and was making a living from odd jobs. Steinhaus later said that Banach was his greatest
discovery. Banach's doctoral thesis (1920, published in 1922) founded the theory of complete normed spaces, now called
Banach spaces, and contained the following theorem, perhaps the most used existence theorem in mathematics.


\begin{theorem}[Banach, 1922]\label{thm:banach}
Let $X$ be a non-empty complete metric space and $f:X\to X$ a $k$-contraction, $k<1$. Then $f$ has a unique fixed point $x^*$. For every
$x_0\in X$ the iterates $x_{n+1}=f(x_n)$ converge to $x^*$, with
\[
d(x_n,x^*)\le\frac{k^n}{1-k}\,d(x_0,x_1)\qquad\text{and}\qquad d(x_n,x^*)\le\frac k{1-k}\,d(x_{n-1},x_n).
\]
\end{theorem}
\begin{proof}
$d(x_{j},x_{j+1})\le k^jd(x_0,x_1)$ by induction, so for $m>n$,
$d(x_n,x_m)\le\sum_{j=n}^{m-1}k^jd(x_0,x_1)\le\frac{k^n}{1-k}d(x_0,x_1)$. The sequence is Cauchy, hence converges to some $x^*$; since $f$ is continuous,
$x^*=\lim x_{n+1}=\lim f(x_n)=f(x^*)$. Letting $m\to\infty$ gives the first bound, and applying it with $x_{n-1}$ as starting point gives the second.
Two fixed points satisfy $d(x^*,y^*)=d(f(x^*),f(y^*))\le k\,d(x^*,y^*)$, so they coincide.
\end{proof}

\begin{think}
Is completeness needed? Is $k<1$ needed, or is $d(f(x),f(y))<d(x,y)$ for $x\ne y$ enough? Try $f(x)=x+\frac1x$ on $[1,+\infty[$.
\end{think}

\begin{example}\label{ex:cosfixed}
$\cos$ maps $[0,1]$ into $[\cos1,1]\subseteq[0,1]$, and $|\cos x-\cos y|\le\sin1\,|x-y|$ there by the mean value theorem (Chapter~4), with
$\sin1<0.85$. So $x=\cos x$ has a unique solution in $[0,1]$, the \emph{Dottie number} $0.7390851332\ldots$, and the a posteriori bound
says when to stop iterating. Convergence is linear, with ratio about $|\sin x^*|\approx0.67$: one digit every six steps. Newton's
method, $x\mapsto x-\frac{x-\cos x}{1+\sin x}$, is a contraction with ratio tending to $0$ near $x^*$, and doubles the number of digits at each
step (Exercise~\ref{ex:newton}). Figure~\ref{fig:cobweb} draws both.
\end{example}

\begin{figure}[ht]
\centering
\fig{ch02-cobweb}
\caption{Iterating $x\mapsto\cos x$ from $x_0=0.1$: the cobweb spirals into the fixed point, the intersection of the graph with the diagonal.
The distances to $x^*$ shrink by a factor close to $|\cos'(x^*)|=\sin x^*\approx0.67$ at each step.}
\label{fig:cobweb}
\end{figure}

\section{Exponential, cosine, sine and \texorpdfstring{$\pi$}{pi}}

A confession is in order. Since Tome~0 we have used $e^{i\theta}$, $\cos$, $\sin$ and $\pi$ as if their properties were
obvious. They are not, and the definition from school, ``$\pi$ is the ratio of a circle's circumference to its
diameter'', presupposes a notion of length of a curve that we have never defined. Here is a cleaner road: define $\exp$
by its power series, let $\cos$ and $\sin$ be its real and imaginary parts on the imaginary axis, and \emph{define} $\pi$
as twice the first positive zero of $\cos$. Everything else, the periodicity, $\cos^2+\sin^2=1$, and even the fact that
the unit circle has length $2\pi$, follows.

Tome~0 admitted the basic properties of $\exp$, $\ln$, $\cos$, $\sin$ and $\pi$. Here they are defined and proved, in the Banach space $\C$.

\begin{lemma}[Cauchy product]\label{lem:cauchyproduct}
If $\sum a_n$ and $\sum b_n$ converge absolutely in $\C$, then $c_n=\sum_{k=0}^na_kb_{n-k}$ defines an absolutely convergent series with
$\sum c_n=\bigl(\sum a_n\bigr)\bigl(\sum b_n\bigr)$.
\end{lemma}
\begin{proof}
$\sum_{n\le N}|c_n|\le\sum_{k,l\le N}|a_k||b_l|\le\sum|a_k|\sum|b_l|$, so $\sum c_n$ converges absolutely. Moreover
$\bigl(\sum_{k\le N}a_k\bigr)\bigl(\sum_{l\le N}b_l\bigr)-\sum_{n\le N}c_n=\sum a_kb_l$ over $k,l\le N$ with $k+l>N$, where $k$ or $l$ exceeds $\frac N2$; so its modulus is at
most $\sum_{k>N/2}|a_k|\sum_l|b_l|+\sum_k|a_k|\sum_{l>N/2}|b_l|\to0$.
\end{proof}

\begin{definition}
For $z\in\C$, $\exp z=\sum_{n\ge0}\frac{z^n}{n!}$, which converges absolutely (compare with $\sum\frac{|z|^n}{n!}$, whose terms are eventually less than
$2^{-n}$). For real $t$, $\cos t=\operatorname{Re}\exp(it)$ and $\sin t=\operatorname{Im}\exp(it)$, that is $\cos t=\sum\frac{(-1)^kt^{2k}}{(2k)!}$ and $\sin t=\sum\frac{(-1)^kt^{2k+1}}{(2k+1)!}$.
We write $e^z=\exp z$ and $e=\exp1$.
\end{definition}

\begin{theorem}\label{thm:exp}
\begin{enumerate}[label=(\alph*)]
\item $\exp(z+w)=\exp z\exp w$, $\exp z\ne0$, $\overline{\exp z}=\exp\bar z$, and $|\exp(it)|=1$ for real $t$, so $\cos^2t+\sin^2t=1$.
\item $\exp$ is differentiable on $\C$ with $\exp'=\exp$; hence $\cos'=-\sin$ and $\sin'=\cos$ on $\R$.
\item On $\R$, $\exp$ is a strictly increasing bijection onto $]0,+\infty[$. Its inverse $\ln$ satisfies $\ln(xy)=\ln x+\ln y$ and $\ln'(x)=\frac1x$.
\item There is a smallest $t>0$ with $\cos t=0$; call it $\frac\pi2$. Then $e^{i\pi/2}=i$, $e^{i\pi}=-1$, $e^{2i\pi}=1$, and $\exp$ is periodic of period $2i\pi$.
\item $e^{it}=1$ for real $t$ iff $t\in2\pi\Z$, and $t\mapsto e^{it}$ maps $[0,2\pi[$ bijectively onto the unit circle.
\item The curve $t\mapsto e^{it}$, $0\le t\le\theta$, has length $\theta$. So $\pi$ is half the perimeter of the unit circle, as in school.
\end{enumerate}
\end{theorem}
\begin{proof}
(a) By Lemma~\ref{lem:cauchyproduct} and the binomial formula, $\exp z\exp w=\sum_n\sum_{k\le n}\frac{z^kw^{n-k}}{k!(n-k)!}=\sum_n\frac{(z+w)^n}{n!}$. So $\exp z\exp(-z)=\exp0=1$.
Conjugation is continuous and commutes with partial sums. Then $|e^{it}|^2=e^{it}\overline{e^{it}}=e^{it}e^{-it}=1$.

(b) By (a),
\[
\begin{aligned}\exp(z+h)-\exp z-h\exp z&=\exp z\,(\exp h-1-h),\\|\exp h-1-h|&\le\sum_{n\ge2}\frac{|h|^n}{n!}\le|h|^2e^{|h|},\end{aligned}
\]
which is $o(|h|)$. For real $t$,
$\frac{\dd}{\dd t}e^{it}=ie^{it}$ by the chain rule for a complex function of a real variable (apply the estimate with $h=is$); take real and
imaginary parts.

(c) For $x>0$, $\exp x>1+x>1$; for $x<y$, $\exp y=\exp x\exp(y-x)>\exp x$, using $\exp x=(\exp\frac x2)^2>0$. As $\exp x\ge1+x$ and $\exp(-x)=\frac1{\exp x}$, $\exp$
tends to $+\infty$ and $0$ at $\pm\infty$. It is continuous and strictly increasing, so by Lemma~\ref{lem:monoivt} below it takes every value in
$]0,+\infty[$. The inverse $\ln$ turns products into sums because $\exp$ turns sums into products, and it is differentiable with
$\ln'(x)=\frac1{\exp'(\ln x)}=\frac1x$ (the derivative of an inverse function, Tome~0).

(d) From the series, $\sin t\ge t-\frac{t^3}6>0$ on $]0,2]$ (the terms after $t$ decrease in modulus and alternate in sign, since $\frac{t^2}{(2k+2)(2k+3)}<1$),
so $\cos$ is strictly decreasing on $[0,2]$. And $\cos2\le1-2+\frac{2^4}{24}=-\frac13<0$ by the same alternating bound. By Lemma~\ref{lem:monoivt}, $\cos$
vanishes at exactly one point of $]0,2[$, and $\cos>0$ on $[0,\frac\pi2[$. Then $\sin\frac\pi2=1$ (it is positive and $\sin^2=1-\cos^2$), so $e^{i\pi/2}=i$, and
powers give $e^{i\pi}=-1$ and $e^{2i\pi}=1$; then $\exp(z+2i\pi)=\exp z$ by (a).

(e) Let $0<t<2\pi$ and $e^{it/4}=u+iv$. As $0<\frac t4<\frac\pi2$, $u=\cos\frac t4>0$ and $v=\sin\frac t4>0$ (sine increases on $[0,\frac\pi2]$ from $0$ since its derivative $\cos$ is
positive). Now $e^{it}=(u+iv)^4=u^4-6u^2v^2+v^4+4iuv(u^2-v^2)$. If it equals $1$, then $u^2=v^2=\frac12$, and the real part is
$\frac14-\frac64+\frac14=-1\ne1$. So $e^{it}\ne1$ on $]0,2\pi[$, and by periodicity $e^{it}=1$ iff $t\in2\pi\Z$. Injectivity on $[0,2\pi[$ follows ($e^{is}=e^{it}$ means
$e^{i(s-t)}=1$). Onto: let $|u+iv|=1$. If $u,v\ge0$, $\cos$ decreases continuously from $1$ to $0$ on $[0,\frac\pi2]$, so $u=\cos t$ for some $t$ there
(Lemma~\ref{lem:monoivt}), and then $v=\sin t$ since both are $\ge0$ with the same square. The other quadrants follow by multiplying by $i$,
$-1$ or $-i$, that is, by adding $\frac\pi2$, $\pi$ or $\frac{3\pi}2$.

(f) The length of a $C^1$ curve $\gamma$ is $\int|\gamma'|$ (Tome~4, Chapter~9), and $|\frac{\dd}{\dd t}e^{it}|=|ie^{it}|=1$.
\end{proof}

The proof used one lemma, a first case of the intermediate value theorem; the general theorem is in Chapter~3.

\begin{lemma}\label{lem:monoivt}
Let $f:[a,b]\to\R$ be continuous and strictly monotone. Then $f$ takes every value between $f(a)$ and $f(b)$. The same holds on an open
interval $]a,b[$ ($-\infty\le a<b\le+\infty$) for the values strictly between the limits of $f$ at $a$ and $b$.
\end{lemma}
\begin{proof}
Say $f$ increases, and $f(a)<y<f(b)$. Let $s=\sup\{x\in[a,b]:f(x)<y\}$. If $f(s)<y$, then $s<b$ and continuity gives $f<y$ on some $[s,s+\eta]$, contradicting
the definition of $s$. If $f(s)>y$, then $s>a$ and $f>y$ on some $[s-\eta,s]$; since $f$ is increasing, $f>y$ on $[s-\eta,b]$, so $s-\eta$ is a smaller upper
bound. Hence $f(s)=y$. On an open interval, pick $a'<b'$ inside with $f(a')<y<f(b')$, which exist by definition of the limits.
\end{proof}

\section{Fractals as fixed points}

Here is a game. Draw any shape in the plane, a cat, a blob, a single point. Shrink it by a factor $2$ towards each corner of a
triangle, keep the three copies, and repeat. After ten rounds the cat has disappeared and the Sierpi\'nski triangle stands in
its place, the same triangle whatever you drew first. It is Dottie's calculator again, played with shapes instead of
numbers: one needs a distance between shapes, and the right one is due to Hausdorff.

\begin{definition}
Let $\mathcal F(X)$ be the set of non-empty closed bounded subsets of a metric space $X$. The \emph{Hausdorff distance} between $A,B\in\mathcal F(X)$ is
\[
d_H(A,B)=\max\Bigl(\sup_{a\in A}d(a,B),\ \sup_{b\in B}d(b,A)\Bigr),
\]
the smallest $r$ such that each set lies in the closed $r$-neighbourhood of the other.
\end{definition}

It is a metric on $\mathcal F(X)$: $d_H(A,B)=0$ forces $A\subseteq\bar B=B$ and $B\subseteq A$; the triangle inequality follows from
$d(a,C)\le d(a,b)+d(b,C)$, taking the infimum over $b\in B$ and then the suprema.

\begin{theorem}\label{thm:hausdorffcomplete}
If $X$ is complete, then $(\mathcal F(X),d_H)$ is complete.
\end{theorem}
\begin{proof}
Let $(A_n)$ be Cauchy, and $A=\bigcap_n\overline{\bigcup_{k\ge n}A_k}$, a closed set. Fix $\varepsilon>0$ and choose $N_0<N_1<\cdots$ with $d_H(A_m,A_n)<\varepsilon2^{-j-1}$ for
$m,n\ge N_j$. \emph{Points of $A_n$ are near $A$.} Let $n\ge N_0$ and $x\in A_n$. Put $x_0=x$, $m_0=n$, and for $j\ge1$ let $m_j=\max(n,N_j)$ and choose $x_j\in A_{m_j}$ with
$d(x_{j-1},x_j)<\varepsilon2^{-j}$; this is possible because $m_{j-1},m_j\ge N_{j-1}$. Then $(x_j)$ is Cauchy and converges to some $y$ with $d(x,y)<\varepsilon$, and $y\in A$,
since for each $p$ it is a limit of points $x_j\in A_{m_j}$ with $m_j\ge p$ eventually. So $A\ne\varnothing$ and $\sup_{x\in A_n}d(x,A)\le\varepsilon$.
\emph{Points of $A$ are near $A_n$.} Let $n\ge N_0$ and $a\in A$. As $a\in\overline{\bigcup_{k\ge n}A_k}$, there are $k\ge n$ and $y\in A_k$ with $d(a,y)<\frac\varepsilon2$, and
$d(y,A_n)\le d_H(A_k,A_n)<\frac\varepsilon2$. So $d(a,A_n)<\varepsilon$. Finally $A$ is bounded, being within $\varepsilon$ of the bounded $A_n$. Hence $A\in\mathcal F(X)$ and
$d_H(A_n,A)\le\varepsilon$ for $n\ge N_0$.
\end{proof}

\begin{theorem}[Hutchinson, 1981]\label{thm:hutchinson}
Let $X$ be complete and $f_1,\dots,f_m:X\to X$ contractions with ratios $k_i<1$. There is a unique $K\in\mathcal F(X)$ with
$K=\overline{f_1(K)\cup\cdots\cup f_m(K)}$, the \emph{attractor} of the system, and for every $A\in\mathcal F(X)$ the iterates $F^n(A)$ of
$F(A)=\overline{\bigcup_if_i(A)}$ converge to $K$ in $d_H$.
\end{theorem}
\begin{proof}
$F(A)$ is closed, non-empty and bounded, so $F$ maps $\mathcal F(X)$ to itself. For each $i$, $d(f_i(a),f_i(B))\le k_id(a,B)$, so $d_H(f_i(A),f_i(B))\le k_id_H(A,B)$,
where $d_H$ is extended to bounded sets by the same formula. A point of $\bigcup f_i(A)$ lies in some $f_i(A)$, within $k_id_H(A,B)$ of $f_i(B)$; so
$d_H(\bigcup f_i(A),\bigcup f_i(B))\le\max_ik_i\,d_H(A,B)$. Closures do not change $d_H$, since $d(x,\bar B)=d(x,B)$ and $d(\cdot,B)$ is continuous. So $F$ is a
contraction of the complete space $\mathcal F(X)$, and Banach's theorem applies.
\end{proof}

For $f_1(x)=\frac x3$ and $f_2(x)=\frac{x+2}3$ on $\R$, the attractor is the \emph{Cantor set} $C$ of the numbers of $[0,1]$ with a ternary
expansion using only the digits $0$ and $2$: $f_1$ and $f_2$ prepend a digit $0$ or $2$, so $f_1(C)\cup f_2(C)=C$, and $C$ is closed and bounded
(Exercise~\ref{ex:cantorifs}). Starting from $[0,1]$, $F^n([0,1])$ is the union of the $2^n$ intervals of the $n$-th stage. Starting from a
single point, it is a set of $2^n$ points, which converges to the same $C$. Figure~\ref{fig:sierpinski} shows the same phenomenon in the
plane.

\begin{figure}[ht]
\centering
\fig{ch02-sierpinski}
\caption{Three contractions of ratio $\frac12$ towards the corners of a triangle. Starting from a square (left), the iterates $F(A)$, $F^2(A)$ and
$F^6(A)$ forget the square and converge to the Sierpi\'nski triangle, the attractor of Theorem~\ref{thm:hutchinson}.}
\label{fig:sierpinski}
\end{figure}

\begin{beauty}
One theorem with a three-line proof, Banach's, solves equations of all kinds: numerical ($x=\cos x$), algebraic (Newton's
method, and Hensel's lemma in $\Z_p$, Exercise~\ref{ex:hensel}), differential (Picard--Lindel\"of, Exercise~\ref{ex:picard}), probabilistic
(PageRank and the stationary laws of Markov chains, numerically) and geometric (fractals). Each time, the work is
to find a complete space on which the problem is a contraction. And completeness itself, the property that made $\R$ out
of $\Q$, makes $\Q_p$ out of $\Q$ for another distance: Ostrowski's theorem (Tome~1, Chapter~12) says that $\R$ and the $\Q_p$ are
the only completions of $\Q$ for an absolute value.
\end{beauty}

\section{Exercises}

\begin{exercise}[$\star$]\label{ex:ultra}
Let $d$ be an ultrametric: $d(x,z)\le\max(d(x,y),d(y,z))$. Show that every triangle is isosceles, with the two longest sides equal; that every
point of a ball is a centre of it; that two balls are either disjoint or nested; and that every open ball is also closed.
\end{exercise}

\begin{exercise}[$\star$]\label{ex:newton}
Let $a>0$ and $N(x)=\frac12\bigl(x+\frac ax\bigr)$. Show that $N$ maps $[\sqrt a,+\infty[$ into itself, is a contraction there with ratio $\frac12$, and that the
errors $e_n=x_n-\sqrt a$ satisfy $e_{n+1}=\frac{e_n^2}{2x_n}$. Starting from $x_0=2$ for $a=2$, how many correct digits does $x_5$ have?
\end{exercise}

\begin{exercise}[$\star$]\label{ex:cantorifs}
Show that the Cantor set $C$ is closed and bounded and satisfies $C=f_1(C)\cup f_2(C)$ for $f_1(x)=\frac x3$, $f_2(x)=\frac{x+2}3$. What is the attractor of
$\frac x2$ and $\frac{x+1}2$? Of $\frac x3$, $\frac{x+1}3$ and $\frac{x+2}3$?
\end{exercise}

\begin{exercise}[$\star\star$]\label{ex:L1incomplete}
On $C([0,1])$, let $\|f\|_1=\int_0^1|f|$. Show that this is a norm, and that $C([0,1])$ is not complete for it: find a Cauchy sequence of
continuous functions that converges, in a sense to be made precise, to the indicator of $[\frac12,1]$.
\end{exercise}

\begin{exercise}[$\star\star$]\label{ex:cantorinter}
Show that $X$ is complete iff every decreasing sequence $F_1\supseteq F_2\supseteq\cdots$ of non-empty closed sets with $\operatorname{diam}F_n\to0$ has a non-empty
intersection (Cantor's intersection theorem). Show that the intersection is then a single point, and that the hypothesis on the
diameters cannot be dropped.
\end{exercise}

\begin{exercise}[$\star\star$]\label{ex:picard}
\emph{Picard--Lindel\"of.} Let $f:[0,T]\times\R\to\R$ be continuous and $L$-Lipschitz in its second variable, and $y_0\in\R$. On $C([0,T])$, consider
$\Phi(y)(t)=y_0+\int_0^tf(s,y(s))\dd s$ and the norm $\|y\|_L=\sup_te^{-2Lt}|y(t)|$. Show that $\Phi$ is a $\frac12$-contraction for $\|\cdot\|_L$, and deduce that
$y'=f(t,y)$, $y(0)=y_0$ has a unique solution on $[0,T]$. What are the iterates for $y'=y$, $y(0)=1$?
\end{exercise}

\begin{exercise}[$\star\star$]\label{ex:hensel}
In $\Z_5$, let $T(x)=x-4(x^2+1)$. Show that $T$ maps $D=\{x:|x-2|_5\le\frac15\}$ into itself and is a contraction there with ratio $\frac15$, and deduce that
$-1$ has a square root in $\Z_5$ congruent to $2$ modulo $5$. Compute its first three $5$-adic digits.
\end{exercise}

\begin{exercise}[$\star\star\star$]\label{ex:devil}
\emph{The devil's staircase as a fixed point.} On the Banach space $B([0,1])$, define $\Phi(f)(x)=\frac12f(3x)$ for $x\in[0,\frac13]$, $\Phi(f)(x)=\frac12$ for
$x\in\,]\frac13,\frac23[$, and $\Phi(f)(x)=\frac12+\frac12f(3x-2)$ for $x\in[\frac23,1]$. Show that $\Phi$ is a $\frac12$-contraction, that its fixed point $c$ is continuous
and non-decreasing with $c(0)=0$ and $c(1)=1$, and that $c$ is constant on each interval removed from $[0,1]$ to build the Cantor set.
\end{exercise}

\section{Hints}

\begin{hint}{ex:ultra}
If $d(x,y)<d(y,z)$, apply the ultrametric inequality to $d(y,z)$ through $x$. For the last claim, a point at distance $\ge r$ from the centre
has a whole ball around it outside.
\end{hint}
\begin{hint}{ex:newton}
$N(x)-\sqrt a=\frac{(x-\sqrt a)^2}{2x}$, and $N'(x)=\frac12(1-\frac a{x^2})\in[0,\frac12[$.
\end{hint}
\begin{hint}{ex:cantorifs}
$C=\bigcap C_n$ with $C_n$ the union of $2^n$ closed intervals. The other two attractors are intervals.
\end{hint}
\begin{hint}{ex:L1incomplete}
Take $f_n$ equal to $0$ on $[0,\frac12-\frac1n]$, $1$ on $[\frac12,1]$, affine between. A limit $g$ in $C([0,1])$ would satisfy $\int|g-\mathbf 1_{[1/2,1]}|=0$.
\end{hint}
\begin{hint}{ex:cantorinter}
For the converse, given a Cauchy sequence, use $F_n=\overline{\{x_k:k\ge n\}}$. For the counterexample, try $F_n=[n,+\infty[$.
\end{hint}
\begin{hint}{ex:picard}
$|\Phi(y)(t)-\Phi(z)(t)|\le L\int_0^te^{2Ls}\|y-z\|_L\dd s$.
\end{hint}
\begin{hint}{ex:hensel}
On $D$, $x=2+5u$ with $|u|\le1$; compute $T(x)-T(x')=(x-x')(1-4(x+x'))$ and look at $1-4(x+x')$ modulo $5$.
\end{hint}
\begin{hint}{ex:devil}
The three pieces of $\Phi(f)-\Phi(g)$ are bounded by $\frac12\|f-g\|_\infty$. The continuous non-decreasing functions with $f(0)=0$, $f(1)=1$ form a closed set that
$\Phi$ maps into itself.
\end{hint}

\section{Solutions}

\begin{solution}{ex:ultra}
\idea Everything follows from one fact: if $d(x,y)\ne d(y,z)$, then $d(x,z)$ is the larger of the two. Balls then behave like the levels of a
tree.
\writeup Let $d(x,y)<d(y,z)$. Then $d(y,z)\le\max(d(y,x),d(x,z))$, and the maximum cannot be $d(x,y)$, so $d(y,z)\le d(x,z)\le\max(d(x,y),d(y,z))=d(y,z)$:
$d(x,z)=d(y,z)$. So the two longest sides are equal. If $y\in B(x,r)$ and $d(z,y)<r$, then $d(z,x)\le\max(d(z,y),d(y,x))<r$, and symmetrically: $B(y,r)=B(x,r)$.
If two balls $B(x,r)$, $B(x',r')$ with $r\le r'$ share a point $z$, then $B(x,r)=B(z,r)\subseteq B(z,r')=B(x',r')$. Finally if $d(y,x)\ge r$, then for
$d(z,y)<r$ we get $d(z,x)=d(y,x)\ge r$ (as $d(z,y)<d(y,x)$): the ball $B(y,r)$ lies outside $B(x,r)$, which is therefore closed.
\end{solution}

\begin{solution}{ex:newton}
\idea Newton's method is a contraction whose ratio is the derivative of $N$; at the fixed point that derivative vanishes, so the
contraction gets stronger and stronger.
\writeup $N(x)-\sqrt a=\frac{x^2-2x\sqrt a+a}{2x}=\frac{(x-\sqrt a)^2}{2x}\ge0$, which is $e_{n+1}=\frac{e_n^2}{2x_n}$ and shows $N([\sqrt a,+\infty[)\subseteq[\sqrt a,+\infty[$. For $x\ge\sqrt a$,
$0\le N'(x)=\frac12(1-\frac a{x^2})<\frac12$, so by the mean value theorem $N$ is a $\frac12$-contraction on this closed, hence complete, set. For $a=2$, $x_0=2$:
$e_0\approx0.586$, and $e_{n+1}\le\frac{e_n^2}{2\sqrt2}$ gives $e_1\le0.12$, $e_2\le5.2\cdot10^{-3}$, $e_3\le9.6\cdot10^{-6}$, $e_4\le3.3\cdot10^{-11}$, $e_5\le3.8\cdot10^{-22}$: $x_5$ has
at least $21$ correct digits (in fact $24$, since $e_1=\frac32-\sqrt2\approx0.086$ is smaller than the bound). The number of digits doubles at each step, while the ratio $\frac12$ alone would only promise one binary digit.
\end{solution}

\begin{solution}{ex:cantorifs}
\idea Write $C$ through its ternary digits; the two maps prepend a digit. For the other systems, guess the attractor from one iterate of
$[0,1]$, and use uniqueness.
\writeup Let $C_0=[0,1]$ and $C_{n+1}=f_1(C_n)\cup f_2(C_n)$. By induction, $C_n$ is the disjoint union of the $2^n$ closed intervals
$I_{d}=[\sum_{k\le n}d_k3^{-k},\sum_{k\le n}d_k3^{-k}+3^{-n}]$ for $d\in\{0,2\}^n$, and $C_{n+1}\subseteq C_n$. We check that $C=\bigcap_nC_n$. If $x=\sum_kd_k3^{-k}$ with all $d_k\in\{0,2\}$,
then $x\in I_{(d_1,\dots,d_n)}$ for every $n$. Conversely, if $x\in\bigcap C_n$, then for each $n$ it lies in exactly one interval of stage $n$ (they are
disjoint), and the interval of stage $n+1$ containing $x$ lies in the one of stage $n$; so the digits $d_k$ are consistent, and $|x-\sum_{k\le n}d_k3^{-k}|\le3^{-n}$
gives $x=\sum_kd_k3^{-k}$. So $C$ is closed (an intersection of closed sets) and bounded. Since $C_{n+1}=f_1(C_n)\cup f_2(C_n)$ and the
$f_i$ are injective, $f_i(\bigcap C_n)=\bigcap f_i(C_n)$; and $f_1(C_n)\subseteq[0,\frac13]$, $f_2(C_n)\subseteq[\frac23,1]$ are disjoint, so
$f_1(C)\cup f_2(C)=\bigcap_n\bigl(f_1(C_n)\cup f_2(C_n)\bigr)=\bigcap_nC_{n+1}=C$. By uniqueness, $C$ is the attractor.
For $\frac x2$, $\frac{x+1}2$: $[0,1]=[0,\frac12]\cup[\frac12,1]$ is fixed, so it is the attractor (binary expansions use all digits). For the three maps of ratio
$\frac13$: $[0,\frac13]\cup[\frac13,\frac23]\cup[\frac23,1]=[0,1]$ again. Fractal attractors need ``missing pieces''.
\end{solution}

\begin{solution}{ex:L1incomplete}
\idea Continuous functions can approximate a jump in the mean, but the jump is not continuous: the limit falls outside the space.
\writeup Positivity: if $f\ne0$ is continuous, $|f|>\frac{|f(x_0)|}2$ on an interval around $x_0$, so $\int|f|>0$. Homogeneity and the triangle inequality come from
those of $|\cdot|$ and the monotonicity of the integral. For $f_n$ as in the hint and $m>n$, $f_n$ and $f_m$ differ only on $[\frac12-\frac1n,\frac12]$, where both are in
$[0,1]$, so $\|f_n-f_m\|_1\le\frac1n$: the sequence is Cauchy. Suppose $f_n\to g$ in $\|\cdot\|_1$ with $g$ continuous. For fixed $\eta>0$, on $[\frac12,1]$ all $f_n=1$, so
$\int_{1/2}^1|1-g|\le\|f_n-g\|_1\to0$, and $g=1$ there by continuity; on $[0,\frac12-\eta]$, $f_n=0$ for $n>\frac1\eta$, so $g=0$ there. As $\eta$ is arbitrary, $g=0$ on $[0,\frac12[$ and
$g(\frac12)=1$, contradicting continuity. The completion is $L^1([0,1])$ (Chapter~9), where $\mathbf 1_{[1/2,1]}$ lives.
\end{solution}

\begin{solution}{ex:cantorinter}
\idea Shrinking closed sets pin down a Cauchy sequence, and conversely the tails of a Cauchy sequence form shrinking closed sets.
\writeup Let $X$ be complete and $(F_n)$ as stated. Pick $x_n\in F_n$. For $m\ge n$, $x_m\in F_n$, so $d(x_n,x_m)\le\operatorname{diam}F_n\to0$: $(x_n)$ is Cauchy, with a limit $x$.
For each $n$, $x$ is a limit of points $x_m\in F_n$ ($m\ge n$), so $x\in F_n$. Two points of $\bigcap F_n$ are at distance $\le\operatorname{diam}F_n$ for all $n$, so they
coincide. Conversely, let $(x_k)$ be Cauchy and $F_n=\overline{\{x_k:k\ge n\}}$; then $\operatorname{diam}F_n=\sup_{k,l\ge n}d(x_k,x_l)\to0$ (the closure does not change the
diameter), so some $x$ lies in every $F_n$: $d(x,x_k)\le\operatorname{diam}F_n$ for $k\ge n$, and $x_k\to x$. Without the diameters: in $\R$, $F_n=[n,+\infty[$ is a
decreasing sequence of non-empty closed sets with empty intersection.
\end{solution}

\begin{solution}{ex:picard}
\idea The integral equation $y=\Phi(y)$ is equivalent to the differential equation, and the weight $e^{-2Lt}$ is chosen so that the
Lipschitz constant $L$ is beaten on the whole interval, not only on a short one.
\writeup $C([0,T])$ with $\|\cdot\|_L$ is complete: $e^{-2LT}\|y\|_\infty\le\|y\|_L\le\|y\|_\infty$, so it has the same Cauchy sequences as the sup norm. $\Phi(y)$ is
continuous (an integral of a continuous function), and
\[
\begin{aligned}|\Phi(y)(t)-\Phi(z)(t)|&\le\int_0^tL|y(s)-z(s)|\dd s\eqbr\le L\|y-z\|_L\int_0^te^{2Ls}\dd s\eqbr\le\frac{e^{2Lt}}2\|y-z\|_L.\end{aligned}
\]
Multiplying by $e^{-2Lt}$ and taking the supremum, $\|\Phi(y)-\Phi(z)\|_L\le\frac12\|y-z\|_L$. Banach gives a unique fixed point $y$. A fixed point is $C^1$ with
$y'=f(t,y)$ and $y(0)=y_0$, by the fundamental theorem of calculus; conversely any solution satisfies $y=\Phi(y)$ by integration, so the
solution is unique. For $y'=y$, $y_0=1$, starting from $1$ the iterates are $\sum_{k\le n}\frac{t^k}{k!}$: Picard's iteration rebuilds the exponential
series.
\end{solution}

\begin{solution}{ex:hensel}
\idea This is Newton's method for $x^2+1$ with the derivative frozen at the first approximation: $2\cdot2=4$, whose inverse modulo $5$ is $4$.
Freezing the derivative loses quadratic convergence but keeps a contraction, which is all Banach needs.
\writeup $D$ is a closed ball of the complete space $\Z_5$, hence complete. For $x\in D$, $x^2+1\equiv4+1\equiv0\pmod5$, so $|x^2+1|_5\le\frac15$ and $|T(x)-2|_5\le\max(|x-2|_5,|4(x^2+1)|_5)\le\frac15$:
$T(D)\subseteq D$. Next $T(x)-T(x')=(x-x')(1-4(x+x'))$, and $x+x'\equiv4\pmod5$, so $1-4(x+x')\equiv1-16\equiv0\pmod5$: $|T(x)-T(x')|_5\le\frac15|x-x'|_5$. The fixed point
satisfies $4(x^2+1)=0$, hence $x^2=-1$, and $x\equiv2\pmod5$. By the a priori bound, $|T^n(2)-x|_5\le5^{-n-1}$, so $T^n(2)\equiv x\pmod{5^{n+1}}$, and an integer congruent to $T^n(2)$ modulo $5^{n+1}$ may replace it
(its image moves by $\frac15$ of that distance at most). Iterating from $2$: $T(2)=2-20=-18\equiv7\pmod{25}$, so $x\equiv7=2+1\cdot5\pmod{25}$; then
$T(7)=7-200=-193\equiv57\pmod{125}$, and $57=2+1\cdot5+2\cdot25$. Check: $57^2+1=3250=26\cdot125$. The digits start $2,1,2$.
\end{solution}

\begin{solution}{ex:devil}
\idea Self-similarity is a fixed point equation, and on a space of bounded functions with the sup norm, halving the values makes
it a contraction. Properties preserved by $\Phi$ and closed under uniform limits pass to the fixed point.
\writeup On each of the three pieces, $|\Phi(f)(x)-\Phi(g)(x)|$ is $\frac12|f(y)-g(y)|$ for some $y$, or $0$; so $\|\Phi(f)-\Phi(g)\|_\infty\le\frac12\|f-g\|_\infty$, and Banach in the
Banach space $B([0,1])$ gives a unique fixed point $c$. Let $M$ be the set of continuous non-decreasing $f$ with $f(0)=0$, $f(1)=1$. It is closed in $B([0,1])$
(uniform limits keep continuity, monotonicity and the values at $0$, $1$), non-empty ($f(x)=x$), and $\Phi(M)\subseteq M$: on $[0,\frac13]$, $\Phi(f)$ rises from $0$ to
$\frac12f(1)=\frac12$, it equals $\frac12$ in the middle, and rises from $\frac12+\frac12f(0)=\frac12$ to $1$ on $[\frac23,1]$, so the pieces join continuously and monotonically. The
iterates of $f\in M$ stay in $M$ and converge to $c$, so $c\in M$. Finally $c=\frac12$ on $]\frac13,\frac23[$, and by the equation $c(x)=\frac12c(3x)$ the function is constant on
each removed interval of $[0,\frac13]$ that is the image of a removed interval under $x\mapsto\frac x3$, and likewise on $[\frac23,1]$; by induction on the stage, $c$ is
constant on every removed interval. This $c$ is the distribution function of Tome~3, Exercise~5.7, and Chapter~10 shows that $c'=0$ almost
everywhere.
\end{solution}

\section*{Sources and further reading}

\begin{itemize}
\item S.~Banach, \emph{Sur les op\'erations dans les ensembles abstraits et leur application aux \'equations int\'egrales}, Fundamenta
Mathematicae 3 (1922), 133--181. The thesis that introduced Banach spaces and the fixed point theorem.
\item M.~Fr\'echet, \emph{Sur quelques points du calcul fonctionnel}, Rendiconti del Circolo Matematico di Palermo 22 (1906), 1--74. Metric
spaces are born.
\item W.~Rudin, \emph{Principles of Mathematical Analysis}, 3rd ed., McGraw--Hill, 1976, Chapters~2, 3 and~8. Metric spaces in forty pages,
and the definition of $\exp$, $\cos$, $\sin$ and $\pi$ by power series followed here (the proof of Theorem~\ref{thm:exp}(e) is Rudin's).
\item F.~Q.~Gouv\^ea, \emph{$p$-adic Numbers: An Introduction}, 2nd ed., Springer, 1997. $\Q_p$ as a completion, with Hensel and Ostrowski.
\item J.~E.~Hutchinson, \emph{Fractals and self-similarity}, Indiana University Mathematics Journal 30 (1981), 713--747. Theorem~\ref{thm:hutchinson} and
much more.
\item M.~F.~Barnsley, \emph{Fractals Everywhere}, Academic Press, 1988. Iterated function systems, the chaos game and the collage
theorem, with pictures.
\item \emph{For the stories in the introduction.} S.~R.~Kaplan, ``The Dottie number'', \emph{Mathematics Magazine} 80 (2007), 73, where the name and its story appear. The
meeting of Steinhaus and Banach in Krak\'ow is told in R.~Ka\l u\.za, \emph{Through a Reporter's Eyes: The Life of Stefan Banach}, Birkh\"auser, 1996.
Hausdorff's name for metric spaces is in his \emph{Grundz\"uge der Mengenlehre}, 1914.
\end{itemize}
